\section{Validating the Neural Posterior Estimator}
\label{sec:validation}
In this section, we check whether our NPE is trained well and produces sensible and well-calibrated posteriors, and measure the shape noise of these posteriors. We then describe how we use these posteriors to produce shear estimates in Sec.~\ref{sec:shear}.

\subsection{Training and Example Posteriors}
\label{ssec:training-examples}

Fig.~\ref{fig:npe-loss} shows the training and validation loss of our model over the course of training. Both losses decrease quickly in the first few epochs, then more slowly for the rest of training. The validation loss fluctuates in the first half of training, when the learning rate is still high, and settles down as the learning rate decreases. It is lowest at epoch 198 out of 200, which we use for all of our analyses. Over the last 10 epochs, the validation loss only improves by 0.008, so the training has converged. It follows the training loss closely throughout, with no sign of overfitting.

\begin{figure}
    \includegraphics[width=\linewidth]{figs/training_loss.pdf}
    \caption{Training (blue) and validation (red) loss of our model as a function of epoch, i.e. the NLL of the true parameters in Eq.~(\ref{eqn:loss-unweighted}). The loss is computed in the rescaled parameter coordinates the flow is trained in, so only differences in the loss are meaningful. The validation loss is lowest at epoch 198 (circle). Inset: a zoom-in on the last 60 epochs.}
    \label{fig:npe-loss}
\end{figure}

Fig.~\ref{fig:posterior-example} shows an example of the $g_1$--$g_2$ posterior from our NPE, for one galaxy from the Xu3 test set. The full 9D posterior of this galaxy is shown in Appendix~\ref{app:corner}. \XW{Overlay Lauren's MCMC contours for this galaxy once they're ready.}

\begin{figure}
    \includegraphics[width=\linewidth]{figs/posterior_example.pdf}
    \caption{Shear posterior of one galaxy from the Xu3 test set (true $\cos i = 0.61$, $\rho_I = 166$, $\rho_S = 23.9$), from $10^4$ posterior samples reweighted by the TF prior. Contours enclose 68\% and 95\% of the posterior, and the dashed lines mark the true shear.}
    \label{fig:posterior-example}
\end{figure}

Once trained, our model draws $10^4$ posterior samples for a galaxy, including the TF reweighting, in about 0.3 seconds on a single NVIDIA V100 GPU. In comparison, recent MCMC fits of DESI datavectors take $\sim$20 minutes per galaxy on one GPU (Willett et al., in prep.; Sec.~\ref{sec:intro}), so our NPE is faster by more than three orders of magnitude.

\subsection{Posterior Calibration}
\label{ssec:calibration}

To test whether our NPE posteriors are well calibrated, we perform a coverage test. If the posteriors are correct, a credible interval at level $L$ should contain the true parameter value for a fraction $L$ of the test galaxies, e.g. 68\% of the 68\% credible intervals should contain the truth. Coverage tests are a standard way to validate approximate posteriors in simulation-based inference \citep{TrustCrisis}, and have been used to validate posteriors in weak lensing, including for NPE \citep{FieldLevelWL, NPEShear}.

We visualize the global coverage of our posteriors with a P--P plot. For each test galaxy $i$ and shear component $g_k$, we compute the probability integral transform (PIT) $u_i$, the TF-weighted fraction of its $10^4$ posterior samples that lie below the true $g_k$. For correct posteriors, the true values are distributed like draws from the posteriors, so the PIT values should be uniformly distributed between 0 and 1 \citep{CookGelmanRubin06, SBC}. The P--P plot shows how often the true value falls below the $u$-quantile of its posterior, which should happen for a fraction $u$ of the galaxies, so the curve should follow the diagonal for calibrated posteriors. If the posteriors are overconfident, the true values fall in their tails too often, giving an S-shaped curve, while a bias in the posteriors shifts the curve to one side of the diagonal. Since the $\vcirc$ of our test galaxies are drawn from the same TFR as our TF prior (Sec.~\ref{ssec:test-set}), this test checks our NPE and the TF reweighting together.

Since we have a large number of test galaxies, we expect the P--P plot to have very little scatter, as we indeed see in the top panel of Fig.~\ref{fig:pp-shear}. Nevertheless, it is important to assess if our scatter is within reasonable bounds. We use the Dvoretzky--Kiefer--Wolfowitz (DKW) inequality to set an upper bound for this scatter \citep{DKW56, Massart90}: for $N$ galaxies with exactly uniform PIT values, the P--P curve stays within $\pm\sqrt{\ln(2/0.05)/2N}$ of the diagonal for all $u$ with at least 95\% probability. This defines a band within which our curve should stay to be considered calibrated, which for our $N=10^5$ Xu3 galaxies is $\pm 0.0043$. The bottom panel of Fig.~\ref{fig:pp-shear} shows that both curves stay within the DKW band for all $u$, deviating from the diagonal by at most 0.003. We therefore consider our NPE shear posteriors to be well calibrated globally.

\begin{figure}
    \includegraphics[width=\linewidth]{figs/pp_shear.pdf}
    \caption{P--P plot of the TF-weighted shear posteriors on the Xu3 test set ($10^5$ galaxies). Top: fraction of galaxies whose PIT value $u_i$, the fraction of posterior samples below the true shear, is at most $u$, for $g_1$ (blue) and $g_2$ (red). Calibrated posteriors follow the diagonal. Bottom: difference from the diagonal. The grey band is the 95\% DKW band, within which the curve of well calibrated posteriors stays with at least 95\% probability.}
    \label{fig:pp-shear}
\end{figure}

For the other seven parameters, we check the calibration at the 16th and 84th percentiles of each posterior, which correspond to $u = 0.16$ and $0.84$ in their P--P plots (Appendix~\ref{app:coverage}). All of them agree with the diagonal to within about 1 percentage point, with the largest deviations of 1.2 points for $\cos i$ and 1.1 points for $R_h$ and $F_{\mathrm{H}\alpha}$, so we consider their posteriors to be accurately calibrated as well. Since we marginalize over all parameters other than the shear when measuring shear, it is the calibration of $g_1$ and $g_2$ that matters in the end.

\subsection{Shape Noise}
\label{ssec:shape-noise}

We measure the shape noise of each galaxy following \citep{KLDESI}, as $\sigma_{\epsilon,i} = \sqrt{\sigma_{g_1,i}^2 + \sigma_{g_2,i}^2}$, where $\sigma_{g_k,i}$ is the standard deviation of the TF-weighted posterior of $g_k$. Across the Xu3 test set, the median $\sigma_{\epsilon,i}$ is 0.066, with 68\% of galaxies between 0.040 and 0.080. About a third of the galaxies have $\sigma_{\epsilon,i} > 0.075$, close to the value of $\sqrt{2}\sigma_\pi = 0.082$ that we would get if the posteriors of both components were equal to the prior. Here $\sigma_\pi = 0.2/\sqrt{12} = 0.058$ is the standard deviation of the uniform prior $\mathcal{U}(-0.1, 0.1)$ on each shear component, and the factor of $\sqrt{2}$ comes from adding the two components in quadrature. The data of these galaxies carry little information about the shear, and as we show in Appendix~\ref{app:posterior-means}, they are the main source of the bias in the posterior means.

Fig.~\ref{fig:shape-noise-maps} shows how the shape noise changes with the quality of the data (panel a), and with the inclination and magnitude of the galaxies (panel b). As expected, galaxies with higher SNR have lower shape noise. Below $\rho_I \approx 400$, where 90\% of our galaxies are, the shape noise depends more on $\rho_S$ than on $\rho_I$, dropping from about 0.078 at $\rho_S < 5$ to about 0.060 at $\rho_S > 50$. This is consistent with \citep{KLInferenceI}, who also found that KL shape noise depends on the emission line SNR. Above $\rho_I \approx 400$, the image SNR also starts to lower the shape noise noticeably, down to about 0.045 for the brightest galaxies with $\rho_S > 20$. In panel (b), the shape noise drops strongly as galaxies become more face-on, from about 0.080 for $\cos i < 0.25$ to about 0.038 for $\cos i > 0.75$. Brighter galaxies also have lower shape noise, about 0.052 for $17 < m_r < 18$ compared with 0.070 for $19 < m_r < 20$.

\begin{figure}
    \includegraphics[width=\linewidth]{figs/shape_noise_maps.pdf}
    \caption{Median shape noise $\sigma_\epsilon$ of the Xu3 test set galaxies (a) in bins of image SNR $\rho_I$ and H$\alpha$ SNR $\rho_S$, and (b) in bins of true $\cos i$ and $m_r$. Bins with fewer than 30 galaxies are left blank. In (a), the solid and dashed contours enclose 68\% and 95\% of the galaxies. Note that the two panels have different color scales.}
    \label{fig:shape-noise-maps}
\end{figure}

To compare with the forecast of \citep{KLDESI}, we compute the effective shape noise of the test set, $\sigma_{\mathrm{eff}} = \left(\frac{1}{N}\sum_i \sigma_{\epsilon,i}^{-2}\right)^{-1/2}$, which weights each galaxy by its inverse variance as in a cosmic shear analysis. Since our test galaxies are drawn from the catalog, they already follow the number density of the sample, which \citep{KLDESI} also weight by. We find $\sigma_{\mathrm{eff}} = 0.041$, compared with the 0.047 forecast for sample 3 in \citep{KLDESI}, even though \citep{KLDESI} use four emission lines, H$\alpha$, the [O\,\textsc{ii}] doublet, and [O\,\textsc{iii}]$\lambda\lambda$4960,5008, while we only use H$\alpha$. We expect this difference to come mainly from our narrower prior on the shear. Galaxies with little shear information have posteriors as wide as our prior $\mathcal{U}(-0.1, 0.1)$ and still contribute to $\sigma_{\mathrm{eff}}$, while with the $\mathcal{U}(-0.5, 0.5)$ prior of \citep{KLDESI}, their posteriors would be much wider and carry almost no weight. As in \citep{KLDESI}, we also find lower shape noise for galaxies that are closer to face-on.
